\subsection{Graded prime ideals}

Let \(A = \bigoplus_{i \geqslant 0} A_i\) be a graded commutative ring with positive degrees and \(m\) the graded ideal \(\bigoplus_{i \geqslant 1} A_i\) ; we shall call the two graded ideals \(a = \bigoplus_{i \geqslant 0} a_i\) , \(b = \bigoplus_{i \geqslant 0} b_i\) of A equivalent if there exists an integer \(n_0\) such that \(a_n = b_n\) for \(n \geqslant n_0\) (clearly it is an equivalence relation). A graded ideal is called essential if it is not equivalent to \(m\) .

PROPOSITION 4. Let \(\mathfrak{p} = \bigoplus_{i\geqslant 0}\mathfrak{p}_i\) be a graded ideal of A; for \(\mathfrak{p}\) to be prime, it is necessary and sufficient that, if \(x\in \mathbf{A}_m,y\in \mathbf{A}_n\) satisfy \(x\notin \mathfrak{p}\) and \(y\notin \mathfrak{p}\) , then \(xy\notin \mathfrak{p}\) .

The condition is obviously necessary. Conversely, if it is fulfilled, then in the graded ring \(A/p = \bigoplus_{i \geq 0} A_i/p_i\) the product of two homogeneous elements \(\neq 0\) is \(\neq 0\) and hence A/p is an integral domain (Algebra, Chapter II, §11, no. 4, Proposition 7).

PROPOSITION 5. Let \(\mathfrak{a} = \bigoplus_{i\geqslant 0}\mathfrak{a}_i\) be a graded ideal of A and \(n_0\) an integer \(>0\) . For there to exist a graded prime ideal \(\mathfrak{p} = \bigoplus_{i\geqslant 0}\mathfrak{p}_i\) such that \(\mathfrak{p}_n = \mathfrak{a}_n\) for \(n\geqslant n_0\) , it is necessary and sufficient that, for pair of homogeneous elements \(x,y\) of degrees \(\geqslant n_0\) , the relation \(xy\in \mathfrak{a}\) implies '' \(x\in \mathfrak{a}\) or \(y\in \mathfrak{a}"\) . If there exists \(n\geqslant n_0\) such that \(\mathfrak{a}_n\neq \mathbf{A}_n\) , the prime ideal satisfying the above conditions is unique.

The condition of the statement is obviously necessary. If \(a_{n} = A_{n}\) for all \(n \geqslant n_{0}\) , clearly every prime ideal containing m is a solution to the problem; there may therefore be several prime ideals which solve the problem; however, any two of these ideals are obviously equivalent. Suppose then that there exists a homogeneous element \(a \in A_{d}\) (with \(d \geqslant n_{0}\) ) not belonging to \(a_{d}\) . Let p be the set of \(x \in A\) such that \(ax \in a\) . Clearly p is an ideal of A; as the homogeneous components of ax are the products by a of those of x and a is a graded ideal, p is a graded ideal; moreover, \(1 \notin p\) and hence \(p \neq A\) . To prove that p is prime, it is sufficient to show that, if \(x \in A_{m}\) , \(y \in A_{n}\) satisfy \(x \notin p\) and \(y \notin p\) , then \(xy \notin p\) (Proposition 4). Then \(ax \notin a_{m+d}\) , \(ay \notin a_{n+d}\) , whence by hypothesis

\[
a ^ {2d} x y \notin \mathfrak {a} _ {m + n + 2 d};
\]

we conclude that \(axy \notin \mathfrak{a}_{m+n+d}\) since \(xy \notin \mathfrak{p}\) . Finally, if \(n \geqslant n_0\) and \(x \in \mathbf{A}_n\) , the conditions \(x \in \mathfrak{a}_n\) and \(ax \in \mathfrak{a}_{n+d}\) are equivalent by hypothesis and hence \(\mathfrak{p} \cap \mathbf{A}_n = \mathfrak{a}_n\) , which completes the proof of the existence of the graded prime ideal \(\mathfrak{p}\) which solves the problem. If also \(\mathfrak{p}'\) is another graded prime ideal of \(\mathbf{A}\) such that \(\mathfrak{p}' \cap \mathbf{A}_n = \mathfrak{a}_n\) for \(n \geqslant n_0\) , then \(a \notin \mathfrak{p}'\) and \(ax \in \mathfrak{p}'\) for all \(x \in \mathfrak{p}\) , whence \(\mathfrak{p} \subset \mathfrak{p}'\) since \(\mathfrak{p}'\) is prime. On the other hand, if \(x\) is a homogeneous element of degree \(n \geqslant 0\) of \(\mathfrak{p}'\) , \(ax\) is homogeneous of degree \(n + d \geqslant n_0\) and therefore belongs to \(\mathfrak{p}' \cap \mathrm{A}_{n+d} = \mathfrak{a}_{n+d}\) , whence by definition \(x \in \mathfrak{p}\) , which shows that \(\mathfrak{p}' \subset \mathfrak{p}\) and finally \(\mathfrak{p}' = \mathfrak{p}\) .

PROPOSITION 6. Let \(d\) be an integer \(\geqslant 1\) .

\begin{enumerate}
\def\labelenumi{(\roman{enumi})}
\item
  For every essential graded ideal \(\mathfrak{p}\) of \(\mathbf{A},\mathfrak{p}\cap \mathbf{A}^{(d)}\) is an essential graded prime ideal of \(\mathbf{A}^{(d)}\) .
\item
  Conversely, for every essential graded prime ideal \(\mathfrak{p}'\) of \(\mathbf{A}^{(d)}\) , there exists a unique (necessarily essential) graded prime ideal \(\mathfrak{p}\) of \(\mathbf{A}\) such that \(\mathfrak{p} \cap \mathbf{A}^{(d)} = \mathfrak{p}'\) .
\item
  If \(a \in \mathbf{A}_k\) does not belong to \(\mathfrak{p}_k\) , \(a^{kd}\) does not belong to \(\mathfrak{p}_{kd}\) and hence \(\mathfrak{p} \cap \mathbf{A}^{(d)}\) is essential.
\item
  For all \(n \geqslant 0\) , the set \(\mathfrak{p} \cap \mathbf{A}_n\) must be equal to the set \(\mathfrak{a}_n\) of \(x \in \mathbf{A}_n\) such that \(x^d \in \mathfrak{p}'\) . Let us show that \(\mathfrak{a} = \bigoplus_{n \geqslant 0} \mathfrak{a}_n\) is a graded prime ideal; as \(\mathfrak{a}_n = \mathfrak{p}_n'\) when \(n\) is a multiple of \(d\) , since \(\mathfrak{p}'\) is prime, this will prove the uniqueness of \(\mathfrak{p}\) . Now, if \(x \in \mathfrak{a}_n, y \in \mathfrak{a}_n, (x - y)^{2d}\) is the sum of terms each of which is a product of \(x^d\) or \(y^d\) by a homogeneous element of degree \(nd\) and hence \((x - y)^{2d} \in \mathfrak{p}'\) and, since \(\mathfrak{p}'\) is prime, \((x - y)^d \in \mathfrak{p}'\) and therefore \(\mathfrak{a}_n\) is a subgroup of \(\mathbf{A}\) . As \(\mathfrak{p}'\) is an ideal of \(\mathbf{A}^{(d)}\) , \(\mathfrak{a}\) is a graded ideal of \(\mathbf{A}\) ; finally, the relation \((xy)^d \in \mathfrak{p}'\) implies \(x^d \in \mathfrak{p}'\) or \(y^d \in \mathfrak{p}'\) , which completes the proof by virtue of Proposition 4.
\end{enumerate}

Let A be a graded commutative ring with positive degrees and p an essential graded prime ideal of A. The set S of homogeneous elements of A not belonging to p is multiplicative and the ring of fractions \(S^{-1}A\) is therefore graded canonically (Chapter II, §2, no. 9) (note that there will in general be homogeneous elements \(\neq0\) of negative degree in this graduation). We shall denote by \(\mathrm{A}_{(\mathrm{p})}\) the subring of \(S^{-1}A\) consisting of the homogeneous elements of degree 0, in other words the set of fractions x/s, where x and s are homogeneous of the same degree in A and \(s\notin p\) . Similarly, for every graded A-module M, \(S^{-1}M\) is graded canonically (loc. cit.) and we shall denote by \(\mathrm{M}_{(\mathrm{p})}\) the subgroup of homogeneous elements of degree 0, which is obviously an \(\mathrm{A}_{(\mathrm{p})}\) -module.

PROPOSITION 7. Let \(\mathfrak{p}\) be a graded prime ideal of \(\mathbf{A}\) , \(d\) an integer \(\geqslant 1\) and \(\mathfrak{p}'\) the graded prime ideal \(\mathfrak{p} \cap \mathrm{A}^{(d)}\) of \(\mathrm{A}^{(d)}\) ; for every graded A-module \(\mathbf{M}\) , the homomorphism \((\mathbf{M}^{(d)})_{(\mathfrak{p}')} \to \mathbf{M}_{(\mathfrak{p})}\) derived from the canonical injection \(\mathbf{M}^{(d)} \to \mathbf{M}\) is bijective.

If S is the set of homogeneous elements of A not belonging to p and \(\mathrm{S}^{\prime}=\mathrm{S}\cap\mathrm{A}^{(d)}\) , the canonical homomorphism \(\phi:\mathrm{S}^{\prime-1}\mathrm{M}^{(d)}\to\mathrm{S}^{-1}\mathrm{M}\) is a homogeneous homomorphism of degree 0 and it is injective, for, if \(x\in M_{nd}\) satisfies sx=0 for \(s\in A_{m}\) , \(s\notin p\) , then also \(s^{d}x=0\) and \(s^{d}\in A_{md}\) , \(s^{d}\notin p^{\prime}\) . It remains to show that the image under \(\phi\) of \((\mathrm{M}^{(d)})_{(\mathfrak{p}^{\prime})}\) is the whole of \(\mathrm{M}_{(\mathfrak{p})}\) ; but if \(x\in M_{n}\) , \(s\in A_{n}\) and \(s\notin p\) , then also \(x/s=(xs^{d-1})/s^{d}\) where \(xs^{d-1}\in A_{nd}\) , \(s^{d}\in A_{nd}\) and \(s^{d}\notin p^{\prime}\) , whence our assertion.

PROPOSITION 8. Let \(\mathfrak{m} = \bigoplus_{i\geqslant 1}\mathbf{A}_i\) ; let \((\mathfrak{p}^{(k)})_{1\leqslant k\leqslant n}\) be a finite family of graded prime ideals of A and a a graded ideal of A such that \(a \cap m \notin p^{(k)}\) for all k; then there exists a homogeneous element \(z \in a \cap m\) not belonging to any of the \(p^{(k)}\) .

We argue by induction on \(n\) , the proposition being trivial for \(n = 1\) . If there exists an index \(j\) such that \(\mathfrak{a} \cap \mathfrak{m} \cap \mathfrak{p}^{(j)}\) is contained in one of the \(\mathfrak{p}^{(k)}\) of index \(k \neq j\) , it follows from the induction hypothesis that there is a homogeneous element \(z' \in \mathfrak{a} \cap \mathfrak{m}\) not belonging to any of the \(\mathfrak{p}^{(k)}\) for \(k \neq j\) and therefore not belonging to \(\mathfrak{p}^{(j)}\) either and this element solves the problem. Suppose then that for every index \(j\) , \(\mathfrak{a} \cap \mathfrak{m} \cap \mathfrak{p}^{(j)}\) is not contained in any of the \(\mathfrak{p}^{(k)}\) of index \(k \neq j\) ; the induction hypothesis therefore implies the existence of a homogeneous element \(y_j \in \mathfrak{a} \cap \mathfrak{m} \cap \mathfrak{p}^{(j)}\) belonging to none of the \(\mathfrak{p}^{(k)}\) for \(k \neq j\) ; as the \(y_j\) are all of degrees \(\geqslant 1\) , we may assume by replacing them by suitable powers (since the \(\mathfrak{p}^{(k)}\) are prime) that \(y_1\) and \(\prod_{j=2}^{n} y_j\) are of the same degree. Then \(z = y_1 + \prod_{j=2}^{n} y_j\) is homogeneous of degree \(\geqslant 1\) and the same argument as in Chapter II, § 1, no. 1, Proposition 2 shows that \(z\) solves the problem.

